% Vietnamese translation for OI Public Library
% Source-File: IOI中国国家候选队论文/2009/罗穗骞/后缀数组——处理字符串的有力工具.ppt
\documentclass[11pt]{article}
\usepackage{oipl-vn}

\newcommand{\Oh}{\mathcal{O}}

\title{Mảng hậu tố: công cụ mạnh để xử lý xâu\\{\large Ghi chú theo slide}}
\author{Luo Suiqian\\Trường Trung học Phụ thuộc Đại học Sư phạm Hoa Nam}
\date{Luận văn đội tuyển dự tuyển quốc gia IOI 2009}

\begin{document}
\maketitle

\origtitle{Suffix arrays as a powerful string tool}
\sourcepath{IOI-China-Candidate-Team-Papers/2009/Luo-Suiqian/suffix-array-slides.ppt}

\section{Nội dung chính}

Slide gồm hai phần lớn: cách xây dựng mảng hậu tố và các ứng dụng dựa trên
mảng \(sa\), hạng \(rank\) và mảng \(height\). Với thi lập trình, bộ ba này
thường thay thế cây hậu tố nhờ cài đặt ngắn hơn và bộ nhớ dễ kiểm soát hơn.

\section{Xây dựng}

Thuật toán nhân đôi sắp xếp hậu tố theo khóa gồm hai nửa độ dài \(2^k\). Sau
mỗi vòng, hạng của hậu tố được cập nhật để dùng cho vòng sau. Thuật toán DC3
có độ phức tạp tuyến tính nhưng cài đặt dài hơn; slide dùng nó để chỉ ra giới
hạn lý thuyết và so sánh với nhân đôi trong thực tế.

\section{Ứng dụng}

\begin{itemize}
  \item Tiền tố chung dài nhất của hai hậu tố được quy về truy vấn nhỏ nhất
  trên đoạn của mảng \(height\).
  \item Bài toán trên một xâu thường đếm hoặc tìm các đoạn lặp.
  \item Bài toán trên hai hay nhiều xâu dùng ký tự phân cách, sau đó xét vị
  trí nguồn của các hậu tố trong một cửa sổ trên mảng \(sa\).
\end{itemize}

\section{Ghi chú}

Khi đọc slide, nên giữ rõ sự khác nhau giữa thứ tự vị trí trong xâu và thứ tự
từ điển trong mảng hậu tố. Nhiều ứng dụng chỉ là biến đổi câu hỏi ban đầu
thành một câu hỏi về các hậu tố kề nhau theo thứ tự từ điển.

\end{document}
